\section{Problemas encontrados}

Durante os testes de inspeção e engenharia dos dados, foram identificados e
mitigados os seguintes problemas:

\begin{enumerate}
  \item \textbf{Inconsistência de codificação (\emph{encoding}) e separadores:}
    Os microdados históricos disponibilizados pelo INEP frequentemente alternam
    entre codificação ISO-8859-1 (Latin-1) e UTF-8, além de utilizarem ponto e
    vírgula (\texttt{;}) como delimitador padrão e vírgula como separador decimal.
    \emph{Solução:} O módulo de leitura foi parametrizado com tolerância a
    erros (\texttt{errors="replace"}) e conversão explícita de tipos numéricos
    com substituição de vírgula por ponto antes da conversão para ponto flutuante.
  \item \textbf{Registros com valores ausentes (\emph{missing values}):}
    Candidatos ausentes nas provas de domingo não possuem notas preenchidas nos
    arquivos brutos, gerando campos vazios que poderiam ocasionar exceções de
    \emph{parsing} ou distorcer o cálculo de médias aritméticas caso fossem
    computados como zero.
    \emph{Solução:} Implementação de filtro obrigatório baseado nos sinalizadores
    \texttt{TP\_PRESENCA\_MT == 1} e \texttt{TP\_PRESENCA\_LC == 1}. Apenas
    candidatos efetivamente presentes entram no acumulador de notas e no
    contador de participantes.
  \item \textbf{Dependência de banda e instabilidade de servidores públicos:}
    O download de arquivos completos de $2$\,GB a partir dos servidores do INEP
    pode sofrer lentidão ou oscilação de rede, prejudicando os testes rápidos
    durante o desenvolvimento do cluster.
    \emph{Solução:} A construção do gerador sintético determinístico com
    \emph{seed} fixa eliminou totalmente o bloqueio das etapas de
    desenvolvimento, permitindo testar a lógica distribuída imediatamente em
    qualquer máquina sem necessidade de download prévio.
\end{enumerate}
